\subsection{Versioned certificate durability under reference drift}
\paragraph{Motivation and scope.}
PARC certificates are defined relative to a frozen candidate universe, a fixed
score, a one-sided support rule, and a specified reference label version.  When
the reference database changes, the truth definition changes as well.  A release
set certified at version \(t_0\) should therefore not be silently inherited at
version \(t_1\).

\paragraph{Versioned labels.}
Let \(P\) be a frozen finite candidate universe.  Each candidate \(p\in P\)
has labels \(Y_p^{(0)},Y_p^{(1)}\in\{0,1\}\), where \(Y_p^{(0)}=1\) means valid
under reference version \(t_0\), and \(Y_p^{(1)}=1\) means valid under reference
version \(t_1\).  For any release set \(R\subseteq P\),
\[
\mathrm{FTR}_v(R)=
\frac{\sum_{p\in R}{\bf 1}\{Y_p^{(v)}=0\}}{|R|\vee 1},\qquad v\in\{0,1\}.
\]
Define
\[
\delta_R^+ =
\frac{\sum_{p\in R}{\bf 1}\{Y_p^{(0)}=1,Y_p^{(1)}=0\}}{|R|\vee 1},
\qquad
\delta_R^- =
\frac{\sum_{p\in R}{\bf 1}\{Y_p^{(0)}=0,Y_p^{(1)}=1\}}{|R|\vee 1}.
\]

\paragraph{Proposition X.1: exact version-shift accounting.}
For any fixed or data-dependent release set \(R\subseteq P\),
\[
\mathrm{FTR}_1(R)=\mathrm{FTR}_0(R)+\delta_R^+-\delta_R^-.
\]
Consequently, \(\mathrm{FTR}_1(R)\le \mathrm{FTR}_0(R)+\delta_R^+\).

\emph{Proof.}
For each candidate \(p\),
\[
{\bf 1}\{Y_p^{(1)}=0\} =
{\bf 1}\{Y_p^{(0)}=0\}
+{\bf 1}\{Y_p^{(0)}=1,Y_p^{(1)}=0\}
-{\bf 1}\{Y_p^{(0)}=0,Y_p^{(1)}=1\}.
\]
Summing over \(p\in R\) and dividing by \(|R|\vee 1\) gives the equality.
Dropping the non-negative term \(\delta_R^-\) gives the inequality.

\paragraph{Corollary X.2: inherited-certificate decay.}
If a \(t_0\)-version PARC release \(R_0\) satisfies
\(\mathbb E[\mathrm{FTR}_0(R_0)]\le \alpha\), then under \(t_1\),
\[
\mathbb E[\mathrm{FTR}_1(R_0)]\le \alpha+\mathbb E[\delta_{R_0}^+].
\]
If an externally calibrated bound \(\delta_{R_0}^+\le \Delta\) is available,
then \(\mathbb E[\mathrm{FTR}_1(R_0)]\le \alpha+\Delta\).  This is a decay
accounting statement, not prospective \(t_1\) alpha control.

\paragraph{Corollary X.3: approximate validity plus reference drift.}
If \(t_0\)-false-candidate e-values obey \(\mathbb E[E_p]\le 1+\eta\), the
same argument yields
\[
\mathbb E[\mathrm{FTR}_1(R_0)]\le \alpha(1+\eta)+\mathbb E[\delta_{R_0}^+].
\]
This separates calibration/e-value error from reference-drift penalty.

\paragraph{Margin-buffer durability.}
Let \(h_p^{(0)}\) and \(h_p^{(1)}\) be versioned energy above hull values, and
let \(Y_p^{(v)}=1\) correspond to \(h_p^{(v)}\le 0\).  For a \(t_0\)-stable
candidate, define \(M_p^{(0)}=-h_p^{(0)}\ge 0\) and
\(D_p=h_p^{(1)}-h_p^{(0)}\).  A \(t_0\)-stable candidate becomes \(t_1\)-unstable
only if \(D_p>M_p^{(0)}\).  Therefore, if a release rule restricts \(t_0\)-stable
released candidates to margin at least \(m\), any stable-to-unstable flip must
satisfy \(D_p>m\).

\paragraph{Corollary X.4: margin-buffer durability bound.}
If \(R_m\) is \(t_0\)-certified and every \(t_0\)-stable released candidate
has \(M_p^{(0)}\ge m\), then
\[
\delta_{R_m}^+
\le
\frac{\sum_{p\in R_m}{\bf 1}\{Y_p^{(0)}=1,D_p>m\}}{|R_m|\vee1}.
\]
If a historical or external drift model gives
\(\Pr(D_p>m\mid p\in R_m,Y_p^{(0)}=1)\le \pi(m)\), then
\[
\mathbb E[\mathrm{FTR}_1(R_m)]\le \alpha+\pi(m),
\]
or \(\alpha(1+\eta)+\pi(m)\) under approximate e-values.  The empirical
\(\pi(m)\) used in this paper is a design diagnostic, not an absolute future
guarantee.

\paragraph{Theorem X.5: versioned recertification.}
Fix the \(t_1\) reference.  If the candidate universe, scores, budget grid,
compatibility rule, and block construction are frozen; \(t_1\)-version
one-sided reliability holds; the PARC exchangeability assumptions hold; and the
returned set satisfies self-consistency, then
\[
\mathbb E[\mathrm{FTR}_1(R_1)]\le \alpha.
\]
If no non-empty compatible set satisfies self-consistency, PARC returns
\(R_1=\varnothing\), for which \(\mathrm{FTR}_1(R_1)=0\).  Recertification
therefore restores a \(t_1\)-relative certificate or refuses to renew the old
release; it does not guarantee a non-empty release.
